\centerline{\Huge Условия и решения задач.}
\centerline{\bf \Huge 8 класс}

\hspace{3mm}


\problem{8.1}Сравните числа X и Y, если $Y = 1 $, а $X = \dfrac{20222022^2+20232023^2}{40454045^2} + \dfrac{1}{\dfrac{40454045^2}{40444044\cdot20232023}}$

\hspace{3mm}

Ответ: равны.

Решение:\\
Пусть $20222022 = a, 20232023 = b,$ тогда дроби переписываются в следующем в виде:\\

$\dfrac{a^2 + b^2}{(a+b)^2} + \dfrac{1}{\dfrac{(a+b)^2}{2ab}} = \dfrac{a^2 + b^2}{(a+b)^2} + \dfrac{2ab}{(a+b)^2} =\dfrac{a^2 + 2ab + b^2}{(a+b)^2}= \dfrac{(a+b)^2}{(a+b)^2} = 1 $. \\

Значит они равны.

\hspace{3mm}

Критерии:\\
Правильное решение через квадрат суммы: \textbf{7 баллов.}\\
Правильное решение подсчетом: \textbf{7 баллов.}\\
Неверное решение или идеи, не продвигающие ход решения: \textbf{0 баллов.}\\
Верный или неверный ответ без обоснования: \textbf{0 баллов.
}

\hspace{3mm}

\problem{8.2}Пять футбольных команд участвуют в круговом турнире. Могут ли команды по завершению турнира набрать следующие очки: 10, 9, 6, 4, 3? (В круговом футбольном турнире каждая команда играет с каждой по одному матчу. За победу дается 3 очка, за ничью — 1, за поражение — 0).

\hspace{3mm}

Решение:\\
Всего было сыграно $\dfrac{5\cdot4}{2} = 10$ матчей. Пусть команда А играет с командой Б. Рассмотрим всевозможные исходы данного матча:

А выиграла - получила 3 очка, Б проиграла - получила 0 очков $\Rightarrow$ Суммарное количество очков всех команд увеличилось на 3.

Б выиграла - получила 3 очка, А проиграла - получила 0 очков $\Rightarrow$ Суммарное количество очков всех команд увеличилось на 3.

А и Б сыграли вничью и получили по 1 очку $\Rightarrow$ Суммарное количество очков всех команд увеличилось на 2.

\newpage

\hspace{3mm}

\noindentЗаметим, что суммарное количество очков всех команд максимально может увеличиваться на 3 за один матч. Значит, за 10 матчей максимальное суммарное количество очков всех команд равняется $10\cdot3=30$. 
Но суммарное количество очков, которые набрали команды: 10 + 9 + 6 + 4 + 3 = 32 > 30, чего не может быть.

\hspace{3mm}

Критерии:\\
Любое правильное решение с верным обоснованием: \textbf{7 баллов.}\\
Доказано, что максимальное суммарное количество очков всех команд равняется 30: \textbf{5 баллов.}\\
Идеи, продвигающие ход решения, но недостаточные для получения верного ответа: \textbf{2 балла.}\\
Посчитано, что всего было сыграно 10 матчей: \textbf{1 балл.}\\
\hspace{3mm}

\problem{8.3}Дан равносторонний треугольник $ABC$. Отрезок $AL$ является продолжением отрезка $AB$ за точку $B$ так, что $AB = BL$. Аналогично $BM$ - продолжение $BC$ за точку $C$ и $CK$ - продолжение отрезка $СА$ за точку $A$ так, что $BC = CM$ и $CA = AK$. Найдите площадь треугольника $KLM$, если известно, что площадь треугольника $ABC$ равна 289.

\hspace{3mm}

Ответ: 2023

Решение:\\


\begin{figure}[h] 
    \includegraphics[scale=0.4]{Городская Олимпиада Элиста 2023/Решения/picture_8.3_sol_2.jpg}
\end{figure}

\newpage

\hspace{3mm}

\noindent\textit{Полезная лемма:} Медиана треугольника делит его на два равновеликих (равных по площади).\\
Доказательство леммы: Пусть нам дан треугольник DEF; медиана $DQ$, проведенная к стороне $EF$, треугольника $DEF$; высота $DH$, проведенная к сторону $EF$, треугольника DEF. Тогда имеем:
$S_{DEQ}=\dfrac{1}{2}\cdot DH\cdot EQ$ и $S_{DFQ}=\dfrac{1}{2}\cdot DH\cdot FQ$, так как $DQ$ - медиана, то $EQ=FQ$, следовательно, $S_{DEQ}=S_{DFQ}$. Лемма доказана.\\

Перейдем к решению задачи, используя полезную лемму:

\begin{figure}[h] 
    \includegraphics[scale=0.4]{Городская Олимпиада Элиста 2023/Решения/picture_8.3_sol.jpg}
\end{figure}

$AB=BL$, следовательно, $CB$ медиана, так как $CB$ - медиана треугольника $ALC$, то $S_{ABC}=S_{CBL}=289$\\
$BC=CM$, следовательно, $LC$ медиана, так как $LC$ - медиана треугольника $LBM$, то $S_{CBL}=S_{CML}=289$\\
$BC=CM$, следовательно, $AC$ медиана, так как $AC$ - медиана треугольника $AMB$, то $S_{ABC}=S_{CAM}=289$\\
$KA=AC$, следовательно, $MA$ медиана, так как $MA$ - медиана треугольника $MCK$, то $S_{CAM}=S_{AKM}=289$\\

\newpage

\hspace{3mm}

\noindent$KA=AC$, следовательно, $BA$ медиана, так как $BA$ - медиана треугольника $KBC$, то $S_{ABC}=S_{KBA}=289$\\
$AB=B$L, следовательно, $KB$ медиана, так как $KB$ - медиана треугольника $KLA$, то $S_{KBA}=S_{KBL}=289$\\
$S_{KLM} = S_{ABC} + S_{CBL} + S_{CML} + S_{CAM} + S_{AKM} + S_{KBA} +S_{KBL} = 289\cdot7 = 2023$


\hspace{3mm}

Критерии:\\
Полностью правильное решение со всеми обоснованиями: \textbf{7 баллов.}\\
Правильное решение без доказательства леммы: \textbf{4 балла.}\\
Идеи, продвигающие ход решения, но недостаточные для получения верного ответа: \textbf{2 балла.}\\
Неверное решение или идеи, не продвигающие ход решения: \textbf{0 баллов.}\\
Верный или неверный ответ без обоснования: \textbf{0 баллов.}

\hspace{3mm}

\problem{8.4}В Элисту приехало 2023 жителя острова ВсОШ, и они сразу же встали в огромный круг друг за другом. Первый сказал: «Следующий за мной Второй является лжецом». Второй сказал: «Следующий за мной Третий является лжецом». Третий сказал: «Следующий за мной Четвёртый является лжецом» и так далее до 2023-го, который сказал: «Следующий за мной Первый является лжецом». Известно, что каждый житель острова - либо рыцарь, который всегда говорит правду, либо лжец, который всегда лжёт, либо хитрец, который может говорить правду или лгать по своему усмотрению. Какое минимальное количество хитрецов могло быть среди 2023-х приезжих жителя острова ВсОШ? 

\hspace{3mm}

Ответ: 1 хитрец.

Решение:\\
\textit{Оценка:} Пусть среди приезжих нету хитрецов. Тогда рассмотрим всевозможные варианты, кем мог быть Первый:\\
1. Предположим, что Первый был рыцарем. Следовательно, Второй был лжецом, а значит, Третий снова был рыцарем, Четвёртый - это лжец и так далее. Каждый нечётный в кругу является рыцарем, а каждый чётный - лжецом. Значит, 2023-ий и Первый оба являются рыцарями, но 2023-ий лжёт, говоря, что Первый является лжецом, но рыцарь не может лгать. Противоречие.


\newpage

\hspace{3mm}

\noindent2. Предположим, что Первый был лжецом. Следовательно, Второй был рыцарем, а значит, Третий снова был лжецом, Четвёртый - это рыцарь и так далее. Каждый нечётный в кругу является лжецом, а каждый чётный - рыцарем. Значит, 2023-ий и Первый оба являются лжецами, но 2023-ий сказал, что Первый является лжецом, то есть он сказал правду, а лжец не может говорить правду. Противоречие.\\
Следовательно, среди приезжих как минимум один хитрец.

\textit{Пример на одного хитреца}: Первый рыцарь, Второй лжец, Третий рыцарь, Четвёртый лжец … 2021-ый рыцарь, 2022-ой лжец, 2023-ий хитрец.

\hspace{3mm}

Критерии:\\
Приведена и доказана правильная оценка на одного хитреца и приведён любой правильный пример с одним хитрецом: \textbf{7 баллов.}\\
Приведена и доказана правильная оценка на одного хитреца, но пример не приведён или приведён неправильный: \textbf{5 баллов.}\\
Приведён любой правильный пример с одним хитрецом, но не доказано, почему не может быть нуля хитрецов: \textbf{2 балла.}\\
Идеи, продвигающие ход решения, но недостаточные для получения верного ответа: \textbf{2 балла.}\\
Приведена оценка не на одного хитреца или приведён пример не с одним хитрецом: \textbf{0 баллов.}\\
Верный или неверный ответ без обоснования: \textbf{0 баллов.}

\hspace{3mm}

\problem{8.5}На международном математическом конгрессе каждый из 200 выступающих математиков рассказал свой доклад ровно один раз 110 другим математикам. Докажите, что хотя бы 2023 пар математиков рассказали доклады друг другу.

\hspace{3mm}

Решение:\\
Количество всевозможных пар математиков равно $\dfrac{200\cdot 199}{2} = 19900.$ Посчитаем, сколько всего было докладов от математиков. Их было ровно $200\cdot110 = 22000.$  Докладов было больше, чем пар, а значит, на какие-то пары пришлось больше одного доклада, но и докладов на одну пару не могло быть более двух (A докладывает B и B докладывает A), а это значит, что у нас будут пары, на которые приходится ровно 

\newpage

\hspace{3mm}

\noindentдва доклада. Следовательно, они и рассказывали доклады друг другу. Таких пар было хотя бы $22000 - 19900 = 2100 > 2023.$

\hspace{3mm}

Критерии:\\
Правильное решение через подсчёт всевозможных пар и общего количества докладов, в котором показано переполнение докладов на пару и указано, что докладов не может быть более двух: \textbf{7 баллов.}\\
Правильное решение через подсчёт всевозможных пар и общего количества докладов, в котором показано переполнение докладов на пару и не указано, что докладов не может быть более двух: \textbf{5 баллов.}\\
Идеи, продвигающие ход решения, но недостаточные для получения верного ответа: \textbf{2 балла.}\\
Неверное решение или идеи, не продвигающие ход решения, такие как "видно, что при таких условиях найдутся 2023 пары" или "очевидно, что найдутся 2023 пары": \textbf{0 баллов.}





